% ============================================================
% manning-style.tex — Stile Manning Publications per libri tecnici
% ============================================================
% Caratteristiche: sans-serif headers, barra laterale codice,
% no-indent paragraphs, numero capitolo grande e grigio.
% ============================================================

% === LINGUA ===
\usepackage[\BabelLang]{babel}
\usepackage[autostyle]{csquotes}

% === BIBLIOGRAFIA ===
\usepackage[
  backend=biber,
  style=numeric,
  sorting=nyt,
  urldate=long,
  maxbibnames=3,
  backref=true,
  backrefstyle=none
]{biblatex}

% === FONT (stile Manning: Times text, Helvetica heads, Inconsolata code) ===
\usepackage{fontspec}
\usepackage{amsmath}
\usepackage{amssymb}
\setmainfont{texgyretermes}[
  Extension=.otf,
  UprightFont=*-regular,
  BoldFont=*-bold,
  ItalicFont=*-italic,
  BoldItalicFont=*-bolditalic
]
\setsansfont{texgyreheros}[
  Extension=.otf,
  Scale=0.92,
  UprightFont=*-regular,
  BoldFont=*-bold,
  ItalicFont=*-italic,
  BoldItalicFont=*-bolditalic
]
\setmonofont{InconsolataN}[
  Extension=.otf,
  Scale=0.88,
  UprightFont=*-Regular,
  BoldFont=*-Bold
]

% === LAYOUT PAGINA (7.375 × 9.25 poll. — formato Manning) ===
\usepackage[
  paperwidth=7.375in,
  paperheight=9.25in,
  top=1in, bottom=1in,
  inner=1.1in, outer=0.85in,
  headheight=14pt,
  footskip=0.5in
]{geometry}

% === SPAZIATURA PARAGRAFI (Manning: no indent, space between) ===
\setlength{\parindent}{0pt}
\setlength{\parskip}{0.5\baselineskip plus 2pt minus 1pt}

% === COLORI (palette Manning — rosso/grigio scuro) ===
\usepackage[dvipsnames,svgnames,table]{xcolor}
\definecolor{manning-red}{HTML}{B71C1C}
\definecolor{manning-dark}{HTML}{212121}
\definecolor{manning-gray}{HTML}{616161}
\definecolor{manning-chap-bg}{HTML}{424242}
\definecolor{code-bg}{HTML}{F4F4F4}
\definecolor{code-bar}{HTML}{BDBDBD}
\definecolor{tip-green}{HTML}{2E7D32}
\definecolor{warn-orange}{HTML}{E65100}
\definecolor{note-blue}{HTML}{1565C0}
\definecolor{link-blue}{HTML}{0D47A1}

% === GRAFICA E TABELLE ===
\usepackage{graphicx}
\graphicspath{{../figures/it/}{../figures/}}
% Inserisce un diagramma; mostra un segnaposto finche' il PNG non e' stato generato
% dalla pipeline dei diagrammi (PlantUML). Uso: \diagramfig{nome-file-senza-estensione}{Didascalia}
\newcommand{\diagramfig}[2]{%
  \begin{figure}[htbp]\centering
  \IfFileExists{../figures/it/#1.png}%
    {\includegraphics[width=0.8\textwidth]{#1.png}}%
    {\fbox{\parbox[c][4cm][c]{0.8\textwidth}{\centering\itshape [Diagramma \texttt{#1.png} --- da generare con il build dei diagrammi]}}}%
  \caption{#2}\label{fig:#1}%
  \end{figure}%
}
\usepackage{float}
\usepackage{booktabs}
\usepackage{longtable}
\usepackage{tabularx}
\usepackage{multirow}
\usepackage{array}
\usepackage{colortbl}

% === CODICE SORGENTE (Manning: sfondo grigio + barra laterale) ===
\usepackage{listings}
\lstset{
  basicstyle=\small\ttfamily\color{manning-dark},
  backgroundcolor=\color{code-bg},
  frame=l,
  framerule=2.5pt,
  rulecolor=\color{code-bar},
  breaklines=true,
  breakatwhitespace=true,
  showstringspaces=false,
  tabsize=2,
  xleftmargin=1em,
  xrightmargin=0.5em,
  aboveskip=0.8em,
  belowskip=0.8em,
  numbers=none,
  numberstyle=\tiny\color{manning-gray},
  numbersep=8pt,
  keepspaces=true,
  keywordstyle=\color{manning-red}\bfseries,
  stringstyle=\color{tip-green},
  commentstyle=\color{manning-gray}\itshape,
  columns=fullflexible,
  literate={€}{{\euro}}1 {→}{{$\rightarrow$}}1 {←}{{$\leftarrow$}}1
}
\lstdefinelanguage{yaml}{
  keywords={true,false,null,yes,no},
  sensitive=false,
  comment=[l]{\#},
  morestring=[b]",
  morestring=[b]',
}

% === INTESTAZIONI E PIÈ DI PAGINA ===
\usepackage{fancyhdr}
\pagestyle{fancy}
\fancyhf{}
\fancyhead[LE]{\small\sffamily\color{manning-gray}\leftmark}
\fancyhead[RO]{\small\sffamily\color{manning-gray}\rightmark}
\fancyfoot[C]{\small\sffamily\thepage}
\renewcommand{\headrulewidth}{0.4pt}
\renewcommand{\headrule}{\hbox to\headwidth{\color{manning-gray}\leaders\hrule height \headrulewidth\hfill}}

% === TITOLI CAPITOLO (Manning: numero grande grigio) ===
\usepackage{titlesec}
\newcommand{\NoTitleHyphenation}{\hyphenpenalty=10000\exhyphenpenalty=10000\relax}
\titleformat{\chapter}[display]
  {\normalfont\sffamily\bfseries\color{manning-dark}\raggedright\NoTitleHyphenation}
  {\fontsize{60}{60}\selectfont\color{manning-gray!40}\thechapter}
  {10pt}
  {\huge}
\titlespacing*{\chapter}{0pt}{0pt}{25pt}
\titleformat{\section}
  {\normalfont\Large\sffamily\bfseries\color{manning-dark}\raggedright\NoTitleHyphenation}
  {\thesection}{0.8em}{}
\titleformat{\subsection}
  {\normalfont\large\sffamily\bfseries\color{manning-dark}\raggedright\NoTitleHyphenation}
  {\thesubsection}{0.8em}{}
\titleformat{\subsubsection}
  {\normalfont\normalsize\sffamily\bfseries\color{manning-gray}\raggedright\NoTitleHyphenation}
  {\thesubsubsection}{0.8em}{}

% === TITOLI PARTI ===
\titleformat{\part}[display]
  {\centering\normalfont\sffamily\bfseries\NoTitleHyphenation}
  {\fontsize{50}{50}\selectfont\color{manning-gray!30}\partname\ \thepart}
  {15pt}
  {\Huge\color{manning-dark}}

% === HYPERLINK ===
\usepackage{xurl}
\usepackage[
  colorlinks=true,
  linkcolor=manning-dark,
  citecolor=manning-red,
  urlcolor=link-blue,
  bookmarks=true,
  bookmarksnumbered=true
]{hyperref}

% === PROTEZIONE OVERFLOW ===
\tolerance=1000
\emergencystretch=1.5em
\setlength{\tabcolsep}{4pt}

% === SPAZIATURA LISTE ===
\usepackage{enumitem}
\setlist{nosep, leftmargin=1.5em, topsep=0.3em, itemsep=0.15em}

% === BOX ADMONITION ===
% ============================================================
% manning-boxes.tex — Box admonition stile Manning
% ============================================================
% Stile: barra laterale colorata, sfondo minimo, icone FontAwesome.
% ============================================================

\providecolor{tip-green}{HTML}{2E7D32}
\providecolor{warn-orange}{HTML}{E65100}
\providecolor{note-blue}{HTML}{1565C0}
\providecolor{manning-gray}{HTML}{616161}

\usepackage{fontawesome5}
\usepackage{tcolorbox}
\tcbuselibrary{skins,breakable}

\newtcolorbox{tipbox}[1][]{
  colback=tip-green!4, colframe=tip-green,
  fonttitle=\sffamily\bfseries,
  title={\faLightbulb\space\textsc{\LabelTip}},
  boxrule=0pt, leftrule=3pt, arc=0pt, breakable,
  coltitle=white, colbacktitle=tip-green,
  toptitle=3pt, bottomtitle=3pt, #1
}

\newtcolorbox{warnbox}[1][]{
  colback=warn-orange!4, colframe=warn-orange,
  fonttitle=\sffamily\bfseries,
  title={\faExclamationTriangle\space\textsc{\LabelWarning}},
  boxrule=0pt, leftrule=3pt, arc=0pt, breakable,
  coltitle=white, colbacktitle=warn-orange,
  toptitle=3pt, bottomtitle=3pt, #1
}

\newtcolorbox{notebox}[1][]{
  colback=note-blue!4, colframe=note-blue,
  fonttitle=\sffamily\bfseries,
  title={\faInfoCircle\space\textsc{\LabelNote}},
  boxrule=0pt, leftrule=3pt, arc=0pt, breakable,
  coltitle=white, colbacktitle=note-blue,
  toptitle=3pt, bottomtitle=3pt, #1
}

\newtcolorbox{dangerbox}[1][]{
  colback=red!3, colframe=red!70!black,
  fonttitle=\sffamily\bfseries,
  title={\faSkullCrossbones\space\textsc{\LabelDanger}},
  boxrule=0pt, leftrule=3pt, arc=0pt, breakable,
  coltitle=white, colbacktitle=red!70!black,
  toptitle=3pt, bottomtitle=3pt, #1
}

\newtcolorbox{versionbox}[1][]{
  colback=purple!3, colframe=purple!60!black,
  fonttitle=\sffamily\bfseries,
  title={\faCogs\space\textsc{\LabelVersionNote}},
  boxrule=0pt, leftrule=3pt, arc=0pt, breakable,
  coltitle=white, colbacktitle=purple!60!black,
  toptitle=3pt, bottomtitle=3pt, #1
}

\newtcolorbox{costbox}[1][]{
  colback=Goldenrod!5, colframe=Goldenrod,
  fonttitle=\sffamily\bfseries,
  title={\faCoins\space\textsc{\LabelCost}},
  boxrule=0pt, leftrule=3pt, arc=0pt, breakable,
  coltitle=white, colbacktitle=Goldenrod,
  toptitle=3pt, bottomtitle=3pt, #1
}

\newtcolorbox{checkpointbox}[1][]{
  colback=tip-green!4, colframe=tip-green,
  fonttitle=\sffamily\bfseries,
  title={\faCheckCircle\space\textsc{\LabelCheckpoint}},
  boxrule=0pt, leftrule=3pt, arc=0pt, breakable,
  coltitle=white, colbacktitle=tip-green,
  toptitle=3pt, bottomtitle=3pt, #1
}

\newtcolorbox{praticabox}[1][]{
  colback=manning-gray!3, colframe=manning-gray,
  fonttitle=\sffamily\bfseries,
  title={\faWrench\space\textsc{\LabelExercise}},
  boxrule=0pt, leftrule=3pt, arc=0pt, breakable,
  coltitle=white, colbacktitle=manning-gray,
  toptitle=3pt, bottomtitle=3pt, #1
}


% === EURO ===
\newcommand{\euro}{€}

% === INDICE ANALITICO ===
\usepackage{makeidx}
\makeindex
